% !TEX root = ../main.tex

\section{Practical Setup}

Table~\ref{tab:defaults} summarizes the main settings used in the current
implementation. These settings define how the tree is built, how evidence is
combined across the tree, and how final cluster decisions are made.

\begin{table}[ht]
\centering
\begin{tabular}{ll}
\toprule
Component & Current choice \\
\midrule
Tree distance & Hamming \\
Linkage & Average linkage \\
Edge significance level & $\alpha_{\mathrm{edge}} = 0.001$ \\
Sibling significance level & $\alpha_{\mathrm{sib}} = 0.01$ \\
Edge test & Projected Wald test with parent-specific whitening \\
Edge dimension rule & Signal count from the parent's local spectrum \\
Sibling test & Projected Wald test with edge-weight adjustment \\
Sibling dimension rule & Smallest positive child signal dimension \\
Sibling basis / calibration & Parent PCA basis, Satterthwaite default \\
Minimum spectral dimension & 2 \\
Include internal spectral rows & On \\
Edge multiple-testing correction & Hierarchical BH correction on the tree \\
Sibling multiple-testing correction & Flat BH on focal pairs \\
Pass-through traversal & On \\
\bottomrule
\end{tabular}
\caption{Main settings used in the current implementation.}
\label{tab:defaults}
\end{table}

Two aspects of the setup deserve special emphasis. First, the two tests do not
combine information in the same way. The child-parent edge test asks whether a
child leaves the background defined by its parent. The sibling test then asks
whether the two children are different enough from one another to justify
keeping the split. Second, once those test tables exist, the final clustering
step is simple: the method walks down the tree and keeps only the splits that
remain supported after correction.

The representation layer should also be read as part of the method. Binary data
enter directly, categorical data enter through one-hot indicator expansion, and
continuous data are discretized before the same testing procedure is applied.
